\documentclass[11pt]{article}
\usepackage[T1]{fontenc}
\usepackage{amsmath,amssymb,amsthm,geometry}
\geometry{margin=.9in}\setlength{\emergencystretch}{2em}
\allowdisplaybreaks[2]
\newtheorem{theorem}{Theorem}\newtheorem{lemma}{Lemma}
\newcommand{\E}{\mathbb E}\newcommand{\dd}{\,d}
\title{Gate transport, overlap, and the weak transverse error}
\author{E5 continuation: analytic partial result}
\date{4 November 2026}
\begin{document}\maketitle

\paragraph{Decision.}
The projected dissipation inequality is not closed through half
response for the actual critical initialization. We remove the
inverse-noise factor from a \emph{signed integrated} gate estimate
and control the overlap and output covariances by a single physical
energy. Smallness of that energy at the moving small-response
hit remains unproved. Moreover, small covariance defects alone do not imply the
desired dissipation inequality: an auxiliary balanced state with
strictly frozen gates shows why a transverse-amplification budget
is also needed. This auxiliary state is not the critical trajectory.
We derive the exact unpaid entry-and-forcing bound on that trajectory.

\section{Actual stopped interval and exact quantities}
Use the critical Gaussian family $S_0=cq^{1/\omega}$, fixed $c>0$,
$\rho\in[13/20,7/10]$, $\lambda=\rho^2$,
$\omega=(1-\lambda)/2$, and the balanced two-layer flow
\[
 f(D)=\int W(U^TD)_+\dd l,\qquad |U|=|W|,\qquad
 S=\int |U|^2\dd l.
\]
The populations have means $e_1,\rho e_2$, covariance $q^2I$,
and mixture weights $1/2$. Write $r=D-f$ and
$\mathcal L=\frac14\sum_p\E_{P_p}|r|^2$.
The four regressions $m,z,h,n$, constants $Q_1=1+q^2$,
$Q=\lambda+q^2$, and parameter gradients $g_m,g_z,g_h,g$
are as in REC/PCD. Thus $h=\E_{P_2}D_2f_1/Q$ and
$n=\E_{P_2}D_2f_2/Q$. Set $w=Q(1-n)$.

Use TBW/RQT's reached interval $I=[\tau_0,\tau]$, where
$n(\tau_0)=\eta_0$, $\mathcal R(\tau_0)\le r_*^2/16$,
$r_*=\lambda/800$, and $\tau$ is the first subsequent hit
of $n=1/2$ or $\mathcal R=r_*^2$. Here
\[
 \mathcal R=4\mathcal L-Q(1-n)^2,
 \quad \tau-t_h\le C_c(1+\ell),\quad
 S\le C_c(1+\ell),\quad \ell=\log(1/q).
\]
Loss decreases with $\mathcal L'=-\|\dot\Theta\|^2$.
No half-response arrival or monotonicity of $n$ is assumed.

For $Q_2=Q$, let
\[
 b_p(U)=Q_p^{-1}\E_{P_p}D_p(U^TD)_+,\qquad
 v_p(U)=Q_p^{-1}\E_{P_p}D_pD\mathbf1_{U^TD>0},
 \quad \theta_p=(v_p)_p.
\]
The current moments of interest are
\[
 C_{12}=\int(U_1)_+(U_2)_+\dd l,\qquad
 T_2=\int W_1W_2|v_2|^2\dd l,\qquad
 O_{12}=\int_{\{U_1>0,U_2>0\}}\theta_2U_1W_2\dd l.
\]
PCD's overlap equation contains $wO_{12}$; its output-moment
equation contains the exact gate term
\begin{equation}\label{an04gtmv1:gate}
 T_2'=A_2+\mathfrak G_2,\quad
 A_2=\int(W_1'W_2+W_1W_2')|v_2|^2\dd l,\quad
 \mathfrak G_2=2\int W_1W_2v_2^TDv_2(U)[U']\dd l.
\end{equation}

\section{First approach: a signed gate budget without an inverse-noise loss}
For any scalar absolutely continuous weight $\chi$ and any
$[a,b]\subset I$, integration by parts in \emph{time} gives exactly
\begin{equation}\label{an04gtmv1:ibp}
 \int_a^b\chi\mathfrak G_2\dd t
 =[\chi T_2]_a^b-\int_a^b\chi'T_2\dd t-\int_a^b\chi A_2\dd t.
\end{equation}
The moving gate has not been discarded: it is accounted for by the
endpoint moment and the two displayed integrals. This operation
does not integrate by parts across a data selector.

REC's feature bound gives $|T_2|\le CS$ and
$|A_2|\le C\sqrt S\|W'\|$. Therefore, with
$D_{ab}=\mathcal L(a)-\mathcal L(b)$,
\begin{align}
 \left|\int_a^b\chi\mathfrak G_2\dd t\right|
 \le C\bigg[&|\chi(a)|S(a)+|\chi(b)|S(b)+\int_a^b|\chi'|S\dd t\nonumber\\
 &+\|\chi\|_\infty\sqrt{D_{ab}\int_a^bS\dd t}\bigg].
 \label{an04gtmv1:signedbudget}
\end{align}
All terms are finite. PCD's explicit Gaussian Hessian and
$|U'|\le C|U|$ justify the chain rule at each positive $q$ by
an integrable bound $CS/q$; positive initialized radii remain
positive at finite times.

Since $D_{ab}\le C$, the actual stopped interval satisfies
\begin{equation}\label{an04gtmv1:gatebounds}
 \left|\int_a^b\mathfrak G_2\dd t\right|\le C_c(1+\ell),
 \qquad
 \left|\int_a^b wh\mathfrak G_2\dd t\right|\le C_c(1+\ell)^2.
\end{equation}
For the second bound, $|wh|\le C$ and
$|(wh)'|\le C\sqrt S\|\dot\Theta\|$, since
$\|g\|,\|g_h\|\le C\sqrt S$ and the regressions are bounded
by loss. Hence $\int S|(wh)'|\le C\sqrt{D_{ab}\int S^3}$.
These estimates remove the former $1/q$ loss, but neither tends
to zero nor pays a fixed residual cap. They bound signed integrals,
not $\int|\mathfrak G_2|$ or its positive part. In particular
they give no pointwise upper-face sign.

\section{Second approach: one energy controls the covariance defects}
Define the physical second-component energy and the hard-gate
weak response by
\[
 E_2=\tfrac12\int(U_2^2+W_2^2)\dd l,\qquad
 n_0=\int W_2(U_2)_+\dd l.
\]
Their difference has an exact nonnegative form
\begin{equation}\label{an04gtmv1:defect}
 \Delta=E_2-n_0
 =\tfrac12\int\{(U_2)_-^2+[W_2-(U_2)_+]^2\}\dd l\ge0,
 \qquad |n-n_0|\le CqS.
\end{equation}
The subscript on $n_0$ denotes a current-state functional, not
the evolution of a substituted zero-noise law.

\begin{theorem}[Covariance control at the actual state]
\label{an04gtmv1:covariances}
Uniformly on the stopped interval,
\begin{equation}\label{an04gtmv1:covbound}
 |C_{12}-z|+|T_2-h|+|O_{12}-z|
                    \le C\sqrt{S\Delta}+CqS.
\end{equation}
Consequently, with $R_\times=Q_1z^2+Qh^2\le\mathcal R$,
\begin{equation}\label{an04gtmv1:source}
 |Q_1zC_{12}+QhT_2-R_\times|
            \le C\sqrt{\mathcal R}(\sqrt{S\Delta}+qS).
\end{equation}
No smallness of $\Delta$ at a moving time is asserted.
\end{theorem}
\begin{proof}
REC proves $b_p=(U_p)_++O(q|U|)$.
Writing $d_U=W_2-(U_2)_+$, one has $\int d_U^2\le2\Delta$,
and
$C_{12}-z=-\int(U_1)_+d_U+\int W_2[(U_1)_+-b_1]$.
Cauchy--Schwarz proves the first estimate.

Let $t=\rho U_2/(q|U|)$. Uniformly in direction,
$\theta_2=\Phi(t)+O(q)$ and $v_2=\theta_2e_2+O(q)$.
For the first identity subtract $1$ from $D_2^2/Q$ inside
the expectation, whose absolute mean is $O(q)$.
The second is REC's exact off-axis moment estimate.
Since $t_+\Phi(-t)$ is bounded, these imply
\[
 (U_2)_+|1-\theta_2|\le Cq|U|,\qquad
 |(U_2)_+|v_2|^2-b_2|\le Cq|U|.
\]
Use
$T_2-h=\int W_1[d_U|v_2|^2+(U_2)_+|v_2|^2-b_2]$.
For $O_{12}-C_{12}$, on $++$ use
$\theta_2W_2-U_2=\theta_2d_U+(\theta_2-1)U_2$.
The same Cauchy--Schwarz estimate pays both. All gate directions,
including $U_2=0$ and either sign of $U_1$, are included.
The last bound follows by multiplying by the signed regressions
and then applying the absolute-value bound, retaining their masses.
\end{proof}

There is also a direct energy balance on the actual flow:
\begin{equation}\label{an04gtmv1:energy}
 E_2'=wn-\tfrac12R_\times+\epsilon_E,
 \qquad |\epsilon_E|\le C\{qS+q^2(1+S)^2\}.
\end{equation}
For completeness the exact expression before estimation is
\[
 E_2'=\E_{\rm mix}[D_2r^T\partial_2f+r_2f_2].
\]
On $P_1$ its first term is $O(qS)$ and its second is
$-Q_1z^2+O(q^2(1+S)^2)$, using the residual projection.
On $P_2$, Euler's identity gives
$D_2r^T\partial_2f=r^Tf-D_1r^T\partial_1f$.
The latter subtraction is $O(qS)$ since $\|\partial_1f\|\le S$
and the residual has bounded Gaussian $L^2$ norm.
The remaining two terms have expectation
$2Qn(1-n)-Qh^2+O(q^2(1+S)^2)$ by orthogonal regression.
Including the mixture halves proves \eqref{an04gtmv1:energy}.

Thus, uniformly for $T\in I$,
\begin{equation}\label{an04gtmv1:clockdefect}
 \Delta(T)=\Delta(\tau_0)+\int_{\tau_0}^T wn\dd t
       -[n(T)-\eta_0]-\tfrac12\int_{\tau_0}^TR_\times\dd t+o(1).
\end{equation}
The error here is bounded by
$C_c[q(1+\ell)^2+q^2(1+\ell)^3]$ from
\eqref{an04gtmv1:defect} and \eqref{an04gtmv1:energy}.
The positive logistic-energy integral cannot be replaced by the
response increase without proving the corresponding clock match.
The existing entrance theorem bounds $\mathcal R(\tau_0)$, not
$\Delta(\tau_0)$. Fixed-resident-time transverse bounds do not
give this moving-time bound. Therefore \eqref{an04gtmv1:covbound}
does not yet supply a small covariance defect throughout $I$.

\section{Why even vanishing defects do not automatically pay dissipation}
The following exact algebra is an audit of a proposed general
implication, \emph{not} a statement about the critical trajectory.
Use the zero-noise two-center model and two effective rank-one
contributions. They can be realized with probability weights
$1/2,1/2$ and vectors multiplied by $\sqrt2$, preserving every
formula below and the gradient-flow time normalization.
Choose
\[
 U_s=(t,-\varepsilon),\quad W_s=(1/t,0),\quad
 t^2=\frac{\sqrt{\varepsilon^4+4}-\varepsilon^2}{2},\qquad
 U_w=W_w=(-a,b).
\]
Both labels are balanced. For $a,b,\varepsilon>0$, each label
activates only its own center, with strict margins. Put
$n=b^2=1/4$ and $a^2=r_*^2/(\lambda n)$.
Then
\[
 m=1,\quad z=0,\quad h=-ab,\quad
 \mathcal R=\lambda h^2=r_*^2,\quad
 C_{12}=O_{12}=0,\quad T_2=h,\quad \Delta=\varepsilon^2/2.
\]
All gates are locally frozen, so their motion term is zero.
The covariance defects tend to zero as $\varepsilon\to0$.
The strong-response projection does not change the weak gradients,
and its completed square is zero. Direct differentiation gives
\begin{equation}\label{an04gtmv1:audit}
 \mathcal D_\perp=\lambda^2h^2(a^2+n),\qquad
 w\mathcal C_\perp=\lambda^2(1-n)h^2,
 \qquad
 \mathcal R'=\lambda^2h^2(1-2n-a^2)>0.
\end{equation}
Indeed $\|g_h\|^2=a^2+n$, $\langle g,g_h\rangle=h$,
and $a^2=\lambda/160000<1/2$. The proposed domination would
require $a^2+2n\ge1$, which fails exactly.

This example refutes only a statewise implication from balance,
small residual, small covariance defects and frozen gates to the
desired dissipation sign. It is not an alternative entrance for
any critical-family theorem. No zero-noise limit is substituted
for the actual positive-noise trajectory. It shows why tracking
the inherited weak cross-output error and its forcing is necessary
in addition to controlling the two gate/overlap terms.

\section{The exact unpaid bound on the critical trajectory}
Retain PCD's projection $P$ away from $g_m$, and set
\[
 A=\|g_m\|^2,\quad B_m=\langle g_m,g_h\rangle,\quad
 H_{ij}=\langle Pg_i,Pg_j\rangle,\quad
 k_h=\langle Pg,Pg_h\rangle.
\]
The denominator $A\ge2m^2/S>0$ on $I$.
Projecting the exact parameter flow gives
\begin{equation}\label{an04gtmv1:hexact}
 h'=a_hh+F_h,\qquad a_h=\tfrac Q2(1-n-H_{hh}),
\end{equation}
where every signed forcing term is retained:
\begin{equation}\label{an04gtmv1:forcing}
 F_h=\frac{B_m}{A}m'-\frac{Q_1z}{2}H_{hz}
        +\frac w2(k_h-h)+\langle Pg_h,PN_{\rm sh}\rangle.
\end{equation}
In deriving this identity, use
$h'=(B_m/A)m'+\langle Pg_h,P\dot\Theta\rangle$ and
$PN=-\tfrac12(Q_1zPg_z+QhPg_h)+PN_{\rm sh}$.
No derivative of a projection and no response monotonicity is used.
The covariance estimate implies
\[
 |k_h-h|\le C\sqrt{S\Delta}+CqS(1+S),
\]
but the new $\Delta$ estimate is not small, and the other signed
forcing terms in \eqref{an04gtmv1:forcing} are not removed.

Let $J(T)=\int_{\tau_0}^Ta_h\dd t$. Exact variation of constants
gives the following \emph{necessary} bound for retention of the cap:
\begin{equation}\label{an04gtmv1:unpaid}
 \boxed{\quad
 \left|h(\tau_0)+\int_{\tau_0}^Te^{-J(t)}F_h(t)\dd t\right|
       \le\frac{r_*}{\sqrt Q}e^{-J(T)}\quad(T\in I).
 }
\end{equation}
This follows from $Qh(T)^2\le\mathcal R(T)\le r_*^2$.
It is an exact obligation to prove in a continuation argument,
not a new estimate derived independently of the stopped cap.
An independent strict improvement is needed to exclude an exit.
For example, replacing $r_*$ on the right by $r_*/2$ would
allocate $r_*^2/4$ to $Qh^2$; the initialized value already
uses at most $r_*^2/16$. This is a concrete sufficient component
budget, not a necessary allocation or a full residual theorem.
It would still leave the other residual components to control.
No estimate in this note establishes such an independent improvement,
or bounds the signed forcing by the shrinking budget when $J$ grows.
Neither a positive $J$ on the actual interval nor its failure is
asserted; the audit explains why positivity cannot be excluded
by the stated statewise constraints.

For the full residual, the aggregate projected face still needs
\[
 w\mathcal C_\perp\le\mathcal D_\perp+\beta^2-\mathcal E_q
 \quad\hbox{at }\mathcal R=r_*^2,\ n<1/2,
\]
with PCD's signed $\mathcal E_q$ of vanishing integrated size.
The gate integration and covariance-energy approach have not
proved this bound; we stop after these two approaches.

\paragraph{Scope and face checks.}
Entry into the residual cap is strict. Its upper face remains
unsigned; the half-response boundary is a stopping target, and a
simultaneous event already reaches that response. Nonnegativity of
$\Delta$ is algebraic, including the face $U_2=0$ and both signs
of $W_2$. The covariance bounds include entry through all input
gates. No invariant quadrant or sign of a mixed moment was assumed.
All constants and noise thresholds may depend on fixed $c,\rho$
and the previous entry choices, with no uniformity as $c\to0$ or
$c\to\infty$. No COD constant is claimed at $s=t_\eta-t_h$.
The new moving-interval bounds follow from exact physical flow,
mass, and loss identities; they do not extend a resident limit to
the weak clock. Half response and a learned probe minimum beating
the coherent comparator remain open. Even a residual-cap proof
would leave the signed probe-motion requirement to be paid.
All prior increments, foundations, checkpoints, and GitHub remain
unchanged. No simulation, empirical value, or numerical certificate
is used as a premise.
\end{document}
